\documentclass[10pt,twocolumn]{article}
\usepackage{neurips_2024}
\usepackage{amsmath,amssymb}
\usepackage{graphicx}
\usepackage{booktabs}
\usepackage{hyperref}
\usepackage{xcolor}
\usepackage{caption}
\usepackage{subcaption}
\usepackage{enumitem}

\setlist{nosep,leftmargin=1.2em}
\captionsetup{font=small,labelfont=bf}

\graphicspath{{../visualizations/}{../results/}}

\title{Conditional Flow Matching vs.\ Diffusion Models:\\A Systematic Comparison on CIFAR-10}

\author{
  Ali Dor \quad Elora Drouilhet \\
  ENS Paris-Saclay / MVA \\
  \texttt{\{ali.dor, elora.drouilhet\}@ens-paris-saclay.fr}
}

\date{}

\begin{document}

\begin{abstract}
We compare Optimal Transport Conditional Flow Matching (OT-CFM) against DDPM and DDIM for image generation on CIFAR-10 using a shared UNet backbone.
OT-CFM achieves a best FID of 10.65 with RK4 at 100 NFE, surpassing DDPM (FID 13.06 at 1000 steps) with 20$\times$ fewer function evaluations.
We extend CFM with class-conditional generation using classifier-free guidance and analyze ODE solver choice, finding midpoint optimal at low NFE budgets and RK4 at high NFE.
Flow matching strictly dominates diffusion at every NFE budget tested.
\end{abstract}

\maketitle

\section{Introduction}

Diffusion models~\cite{ho2020denoising} achieve excellent image quality but require ${\sim}$1000 denoising steps.
DDIM~\cite{song2021denoising} enables fewer-step sampling but quality saturates quickly.
Flow matching~\cite{lipman2023flow,liu2023flow} learns a velocity field transporting noise to data along simple paths.
With optimal transport coupling~\cite{tong2024improving}, OT-CFM produces straight trajectories easily integrated by ODE solvers in far fewer steps.

\textbf{Contributions.}
(1)~We implement OT-CFM from scratch following~\cite{lipman2024flow} and compare against DDPM/DDIM with a shared architecture;
(2)~we extend CFM to class-conditional generation with CFG~\cite{ho2022classifier};
(3)~we systematically analyze ODE solvers across NFE budgets;
(4)~we show flow matching dominates diffusion at every NFE on CIFAR-10.

\section{Background}

\textbf{DDPM.}
The forward process is $q(\mathbf{x}_t|\mathbf{x}_0) = \mathcal{N}(\sqrt{\bar\alpha_t}\mathbf{x}_0, (1{-}\bar\alpha_t)\mathbf{I})$.
A network $\epsilon_\theta$ predicts the noise via $\mathcal{L} = \mathbb{E}\|\epsilon - \epsilon_\theta(\mathbf{x}_t, t)\|^2$, and sampling denoises over $T{=}1000$ steps~\cite{ho2020denoising}.

\textbf{DDIM}~\cite{song2021denoising} reinterprets the reverse process as a deterministic ODE, enabling fewer-step sampling from a DDPM-trained model.

\textbf{CFM.}
Flow matching~\cite{lipman2023flow} learns $v_\theta(\mathbf{x}_t, t)$ generating a path from $p_0$ to $p_1$.
With OT interpolation~\cite{tong2024improving}: $\mathbf{x}_t = (1{-}t)\mathbf{x}_0 + t\mathbf{x}_1$, target $\mathbf{u}_t = \mathbf{x}_1 - \mathbf{x}_0$, loss $\mathcal{L}_{\text{CFM}} = \mathbb{E}\|v_\theta(\mathbf{x}_t, t) - \mathbf{u}_t\|^2$.
Sampling solves the ODE $d\mathbf{x}/dt = v_\theta(\mathbf{x}_t, t)$ from $t{=}0$ to $1$.

\textbf{Solvers.}
Euler (1~NFE/step), Midpoint (2~NFE/step), RK4 (4~NFE/step).

\section{Methods}

\textbf{Architecture.}
All models use a UNet~\cite{ronneberger2015unet}: 128 base channels, multipliers $[1,2,2,2]$, 4 attention heads, GroupNorm-32, sinusoidal time embeddings.
The class-conditional variant uses $[1,2,4]$ multipliers, GroupNorm-8, with learned class embeddings added to the time embedding.

\textbf{Training.}
DDPM and unconditional OT-CFM: 500 epochs, EMA decay 0.9999, Adam with lr $2{\times}10^{-4}$, batch size 128.
Class-conditional CFM: 210 epochs, EMA 0.999, 10\% class dropout for CFG.

\textbf{Evaluation.}
10k generated samples; FID~\cite{heusel2017gans} and IS~\cite{salimans2016improved} vs.\ CIFAR-10 training set; NFE sweep $\{10, 20, 50, 100, 1000\}$.

\section{Results}

\begin{table}[t]
\centering
\caption{Full comparison on CIFAR-10. Best FID per NFE in \textbf{bold}.}
\label{tab:full}
\scriptsize
\setlength{\tabcolsep}{3pt}
\begin{tabular}{@{}llrrr@{}}
\toprule
Method & Solver & NFE & FID\,$\downarrow$ & IS\,$\uparrow$ \\
\midrule
DDPM & DDPM & 1000 & 13.06 & 8.18 \\
\midrule
DDIM & DDIM & 10 & 22.18 & 7.77 \\
DDIM & DDIM & 20 & 18.95 & 7.82 \\
DDIM & DDIM & 50 & 17.96 & 7.84 \\
DDIM & DDIM & 100 & 18.43 & 7.88 \\
\midrule
OT-CFM & Midpoint & 10 & \textbf{17.53} & 8.02 \\
OT-CFM & Midpoint & 20 & \textbf{14.80} & 8.22 \\
OT-CFM & Midpoint & 50 & \textbf{12.02} & 8.44 \\
OT-CFM & Midpoint & 100 & 11.21 & 8.50 \\
OT-CFM & RK4 & 100 & \textbf{10.65} & 8.58 \\
\midrule
Cond.\ CFM & Midpoint & 10 & 17.53 & 8.02 \\
Cond.\ CFM & Midpoint & 50 & 13.17 & 8.28 \\
Cond.\ CFM & RK4 & 100 & 12.19 & 8.46 \\
\bottomrule
\end{tabular}
\end{table}

\begin{table}[t]
\centering
\caption{Best FID per model family.}
\label{tab:best}
\scriptsize
\setlength{\tabcolsep}{3pt}
\begin{tabular}{@{}lrrll@{}}
\toprule
Model & FID & NFE & Solver & ms/img \\
\midrule
DDPM & 13.06 & 1000 & DDPM & 205.0 \\
DDIM & 17.96 & 50 & DDIM & 10.5 \\
OT-CFM & \textbf{10.65} & 100 & RK4 & 20.7 \\
Cond.\ CFM & 12.19 & 100 & RK4 & 24.3 \\
\bottomrule
\end{tabular}
\end{table}

Table~\ref{tab:full} compares all methods. OT-CFM with RK4 at 100 NFE achieves the best FID of \textbf{10.65}. Table~\ref{tab:best} summarizes the best result per model.

\begin{figure}[t]
\centering
\includegraphics[width=\columnwidth]{vector_fields_2d.png}
\caption{Learned OT-CFM velocity fields on a 2D toy problem at different timesteps, showing straight transport paths from noise to data.}
\label{fig:vector_fields}
\end{figure}

\begin{figure}[t]
\centering
\includegraphics[width=\columnwidth]{coupling_comparison.png}
\caption{Independent vs.\ OT coupling. OT produces straighter, non-crossing trajectories that are easier for ODE solvers to integrate.}
\label{fig:coupling}
\end{figure}

\begin{figure}[t]
\centering
\includegraphics[width=\columnwidth]{interpolation_cfm.png}
\caption{Latent interpolation with OT-CFM on CIFAR-10. Smooth transitions confirm a well-structured latent space.}
\label{fig:interp}
\end{figure}

\begin{figure}[t]
\centering
\includegraphics[width=\columnwidth]{conditional_cfg_comparison.png}
\caption{Class-conditional CFM with varying CFG scale. Higher guidance improves class fidelity at the cost of diversity.}
\label{fig:cfg}
\end{figure}

\section{Analysis}

\textbf{CFM dominates diffusion.}
At NFE=100, OT-CFM midpoint achieves FID 11.21 vs.\ DDIM's 18.43 (39\% better).
At just 20 NFE, OT-CFM (14.80) already beats the best DDIM at any NFE (17.96).
DDPM's 1000-step FID of 13.06 is surpassed by OT-CFM at ${\sim}$50 steps---a \textbf{20$\times$ speedup}.

\textbf{Solver choice.}
Midpoint gives the best quality-per-NFE at low budgets.
RK4 achieves the best absolute FID (10.65) at 100 NFE but wastes compute at low NFE (only 5 steps at NFE=20).

\textbf{DDIM saturation.}
DDIM improves to NFE=50 (FID 17.96) then \emph{degrades} at 100 (18.43)---accumulated errors from the stochastic training objective. CFM does not exhibit this.

\textbf{Conditional CFM.}
Slightly worse than unconditional (12.64 vs.\ 11.21 at NFE=100) due to fewer training epochs (210 vs.\ 500). With CFG scale $>1$ the gap would close.

\section{Limitations}

Single dataset (CIFAR-10, 32$\times$32); no confidence intervals (single runs); conditional model undertrained (210 vs.\ 500 epochs); no $\sigma_{\min}$ exploration.

\section{Conclusion}

OT-CFM strictly outperforms DDPM/DDIM at every NFE budget, achieving FID 10.65 at 100 NFE vs.\ DDPM's 13.06 at 1000---a 20$\times$ speedup.
Midpoint is the recommended default solver; class-conditional CFM with CFG provides an orthogonal improvement axis.

\subsection*{Contributions}
\textbf{Ali Dor:} DDPM/DDIM implementation, training, evaluation, all visualizations, results analysis.
\textbf{Elora Drouilhet:} OT-CFM and class-conditional CFM with CFG implementation, training, evaluation.

{\small
\bibliographystyle{plainnat}
\bibliography{references}
}

\onecolumn
\appendix
\section*{NeurIPS Paper Checklist}
\small
\begin{enumerate}[leftmargin=1.5em]
\item \textbf{Claims.} \textit{Yes.} All claims supported by Section~4.
\item \textbf{Limitations.} \textit{Yes.} See Section~6.
\item \textbf{Theory.} \textit{N/A.} No novel theory.
\item \textbf{Experiments.} \textit{Yes.} Full details in Section~3.
\item \textbf{Open access.} \textit{Yes.} Code in project repository; CIFAR-10 is public.
\item \textbf{Experimental details.} \textit{Yes.} Section~3.
\item \textbf{Statistical significance.} \textit{No.} Single runs, noted as limitation.
\item \textbf{Compute resources.} \textit{Yes.} Google Colab T4/A100; runtimes in Table~\ref{tab:best}.
\item \textbf{Code of ethics.} \textit{Yes.}
\item \textbf{Broader impacts.} \textit{N/A.} Course project on standard benchmark.
\item \textbf{Safeguards.} \textit{N/A.} 32$\times$32 models cannot generate realistic content.
\item \textbf{Licenses.} \textit{Yes.} All assets properly cited.
\item \textbf{New assets.} \textit{N/A.}
\item \textbf{Human subjects.} \textit{N/A.}
\item \textbf{IRB.} \textit{N/A.}
\item \textbf{LLM usage.} \textit{Yes.} Claude assisted with code development and formatting.
\end{enumerate}

\end{document}
